\jlauChapterAfterPart{作物表型数据融合研究综述}

\section{研究背景}

作物表型是基因型、环境与管理措施共同作用的外在表现。传统表型调查依赖人工测量，能够获得较高可信度的局部指标，但在大规模田间试验中容易受到时间成本和主观差异限制。随着无人机、地面平台、固定传感器和图像识别方法的发展，研究人员能够更连续地记录冠层结构、长势指数和环境变化。

\section{国内外研究现状}

公开研究显示，多源表型融合通常需要解决三个问题：其一，不同传感器采样时间不一致；其二，图像特征和人工指标存在尺度差异；其三，田间环境会放大缺失值和异常值对模型的影响。常用方法包括标准化处理、时间窗口对齐、特征筛选和集成模型等。

\subsection{传感观测与图像特征}

田间图像可以用于提取冠层覆盖度、颜色指数和纹理特征；环境传感数据则记录温度、湿度、光照和土壤水分。二者结合后能够更完整地描述作物生长状态。

\subsection{数据融合与质量控制}

数据融合的关键不在于简单叠加变量，而在于识别变量之间的互补关系。质量控制阶段应记录缺失机制、异常来源和人工复核规则，避免模型把采集误差误认为真实生物学差异。

\section{小结}

现有研究为智慧农业中的表型数据分析提供了丰富方法，但在跨尺度融合、可解释评估和模板化复现方面仍有改进空间。本示例后续章节围绕这些问题设计可替换的论文骨架。
